\def\level{Première}  % Classe à droite
\def\course{NSI}
\def\chaptername{Listes et Tuples} % Titre au centre
\def\docname{AP05-S}        % Identifiant à gauche

\pagestyle{cours_FP}
\makeheader{} % Affiche le bandeau

\begin{definition}[Les listes]{}

\begin{tabular}{|p{2cm}|p{1cm}|p{1cm}|p{1cm}|p{1cm}|p{1cm}|p{1cm}|p{1cm}|p{1cm}|}
\hline
Indice \verb|i| & 0 & 1 & 2 & 3 & 4 & 5 & 6 & 7 \\ \hline
Elément & \verb|T[0]| & \verb|T[1]| & \verb|T[2]| & \verb|T[3]| & \verb|T[4]| & \verb|T[5]| & \verb|T[6]| & \verb|T[7]| \\ \hline

\end{tabular}



\end{definition}


	\begin{minipage}{0.48\linewidth}
\begin{easybox}{Accéder aux éléments d'une liste}{}
		\begin{CodePythontexAlt}[Largeur=1\linewidth,Centre,PremLigne=1]{}
ma_liste = [1, 2, 3, 4, 5]
print(ma_liste[0]) 	# Affiche: 1
print(ma_liste[2])  # Affiche: 3
print(len(ma_liste))  # Affiche: 5
		\end{CodePythontexAlt}
		\end{easybox}
\end{minipage}\hfill
	\begin{minipage}{0.48\linewidth}
		\begin{easybox}{Modifier les éléments d'une liste}{}
	\begin{CodePythontexAlt}[Largeur=1\linewidth,Centre,PremLigne=1]{}
ma_liste = [1, 2, 3]
ma_liste[0] = 10 
print(ma_liste) 	# Affiche: [10, 2, 3]
print(len(ma_liste))  # Affiche: 3
		\end{CodePythontexAlt}
	\end{easybox}
	
	\end{minipage}

\begin{easybox}{Itérer sur une liste avec une boucle FOR}{}


	\begin{minipage}{0.48\linewidth}
		\begin{CodePythontexAlt}[Largeur=1\linewidth,Centre,PremLigne=1]{}
ma_liste = [4, 8, 7, 9, 4]
for i in range(0, len(ma_liste)):
	print(ma_liste[i])
		\end{CodePythontexAlt}
	\end{minipage}\hfill
	\begin{minipage}{0.48\linewidth}
		\begin{CodePythontexAlt}[Largeur=1\linewidth,Centre,PremLigne=1]{}
ma_liste = [4, 8, 7, 9, 4]
for elt in ma_liste:
	print(elt)
		\end{CodePythontexAlt}
	\end{minipage}
\end{easybox}

\begin{minipage}{0.48\linewidth}
\begin{easybox}[Sommes des éléments d'une liste]{}
		\begin{CodePythontexAlt}[Largeur=1\linewidth,Centre,PremLigne=1]{}
def somme(liste):
	somme = 0
	for elt in liste:
		somme = somme + elt
	return somme
#
		\end{CodePythontexAlt}
\end{easybox}
\end{minipage}\hfill
	\begin{minipage}{0.48\linewidth}
		\begin{easybox}{Recherche de maximum dans une liste}{}
	\begin{CodePythontexAlt}[Largeur=1\linewidth,Centre,PremLigne=1]{}
def maximum(liste):
	maximum = liste[0]
	for elt in liste:
		if elt > maximum:
			maximum = elt
	return maximum
\end{CodePythontexAlt}
	\end{easybox}
	\end{minipage}

	\begin{minipage}{0.48\linewidth}
\begin{easybox}[Recherche présence d'un élément dans une liste]{}
		\begin{CodePythontexAlt}[Largeur=1\linewidth,Centre,PremLigne=1]{}
def est_present(elt, liste):
	for e in liste:
		if e == elt:
			return True
	return False
		\end{CodePythontexAlt}
\end{easybox}
\end{minipage}\hfill
	\begin{minipage}{0.48\linewidth}
\begin{easybox}[Vérification liste triée]{}
		\begin{CodePythontexAlt}[Largeur=1\linewidth,Centre,PremLigne=1]{}
def est_triee(liste):
	for i in range(0, len(liste)-1):
		if liste[i] > liste[i+1]:
			return False
	return True
		\end{CodePythontexAlt}
\end{easybox}
\end{minipage}
			
\begin{easybox}{Liste de listes}{}

		\begin{CodePythontexAlt}[Largeur=1\linewidth,Centre,PremLigne=1]{}
ma_liste = [[1, 2, 3], [4, 5], [7, 8, 9, 10]]
print(len(ma_liste))	# Affiche: 3
print(len(ma_liste[1])) # Affiche: 2
print(ma_liste[0]) 		# Affiche: [1, 2, 3]
print(ma_liste[1])		# Affiche: [4, 5]
print(ma_liste[2])		# Affiche: [7, 8, 9, 10]
print(ma_liste[0][1]) 	# Affiche: 2
print(ma_liste[1][1]) 	# Affiche: 5
print(ma_liste[2][0]) 	# Affiche: 7
		\end{CodePythontexAlt}

		

\end{easybox}

\pagebreak

\begin{minipage}{0.47\linewidth}
\begin{easybox}[Sommes des éléments d'une liste de listes]{}
		\begin{CodePythontexAlt}[Largeur=1\linewidth,Centre,PremLigne=1]{}
def somme(matrice):
	somme = 0
	for liste in matrice:
		for elt in liste:
			somme = somme + elt
	return somme
#
		\end{CodePythontexAlt}
\end{easybox}
\end{minipage}\hfill
	\begin{minipage}{0.51\linewidth}
		\begin{easybox}{Recherche de maximum liste de listes}{}
	\begin{CodePythontexAlt}[Largeur=1\linewidth,Centre,PremLigne=1]{}
def maximum(matrice):
	maximum = matrice[0][0]
	for i in range(0, len(matrice)):
		for j in range(0, len(matrice[i])):
			if matrice[i][j] > maximum:
				maximum = matrice[i][j]
	return maximum
\end{CodePythontexAlt}
	\end{easybox}
	\end{minipage}


\begin{easybox}{Quelques méthodes sur les listes}{}

	\begin{minipage}{0.65\linewidth}
		\begin{itemize}
\item \verb|append(elt)|: ajoute l'élément \verb|elt| à la fin de la liste	
\item \verb|pop()|: supprime et renvoie le dernier élément de la liste
\item \verb|pop(i)|: supprime et renvoie l'élément d'indice \verb|i| de la liste
\end{itemize}

	\end{minipage}\hfill
	\begin{minipage}{0.35\linewidth}
	\begin{CodePythontexAlt}[Largeur=1\linewidth,Centre,PremLigne=1]{}
ma_liste = [1, 2, 3, 4, 5]
ma_liste.append(6)
a = ma_liste.pop()
b = ma_liste.pop(2)
		\end{CodePythontexAlt}
	\end{minipage}
\end{easybox}

\begin{easybox}{Créer une liste par itération}{}

\begin{minipage}{0.48\linewidth}
		\begin{CodePythontexAlt}[Largeur=1\linewidth,Centre,PremLigne=1]{}
ma_liste = []
for i in range(0, 5):
	ma_liste.append(i)

print(ma_liste)		
# Affiche: [0, 1, 2, 3, 4]
		\end{CodePythontexAlt}
	\end{minipage}\hfill
	\begin{minipage}{0.48\linewidth}
		\begin{CodePythontexAlt}[Largeur=1\linewidth,Centre,PremLigne=1]{}
ma_liste = []
for i in range(0, 5):
	ma_liste = [i] + ma_liste

print(ma_liste)
# Affiche: [4, 3, 2, 1, 0]
		
		\end{CodePythontexAlt}
	\end{minipage}
\end{easybox}

\begin{easybox}{Créer une liste par compréhension}{}

	\begin{CodePythontexAlt}[Largeur=1\linewidth,Centre,PremLigne=1]{}
ma_liste = [i for i in range(0, 5)]
print(ma_liste)
# Affiche: [0, 1, 2, 3, 4]
		\end{CodePythontexAlt}
		\end{easybox}


	
\begin{easybox}{Le hasard et les listes}{}
	\begin{itemize}
		\item La fonction \verb|random.choice(liste)| permet de choisir un élément au hasard dans une liste.
		\item La fonction \verb|random.shuffle(liste)| permet de mélanger les éléments d'une liste de manière aléatoire.
	\end{itemize}
\end{easybox}

\begin{easybox}{Le slicing}{}
	Le slicing permet de créer une nouvelle liste à partir d'une liste existante en sélectionnant une portion de cette liste.

	\begin{CodePythontexAlt}[Largeur=1\linewidth,Centre,PremLigne=1]{}
ma_liste = [1, 2, 3, 4, 5]
print(ma_liste[1:4])	# Affiche: [2, 3, 4]
print(ma_liste[:3])		# Affiche: [1, 2, 3]
print(ma_liste[2:4])	# Affiche: [3, 4]
		\end{CodePythontexAlt}
\end{easybox}

		\begin{definition}[Les tuples]{}
		Un \textbf{tuple} est une structure de donnée qui contient des éléments ou items identifiés par des indices.	
		Les tuples ont les mêmes caractéristiques que les listes, à la différence que les éléments d'un tuple ne peuvent pas être modifiés après la création du tuple. Un tuple est une structure de donnée immuable.

			
		\end{definition}
